\documentclass[aspectratio=169,10pt,english]{beamer}

\input{../../../_preambulo_beamer.tex}
\input{../../../_comandos_beamer.tex}

\selectlanguage{english}

\title{MA1001B - Statistical Modeling for Decision Making}
\subtitle{Lesson 24: Central Limit Theorem}
\author{}
\institute{Tecnol\'ogico de Monterrey}
\date{}

\begin{document}

\begin{frame}[plain]
  \vspace{1.2cm}
  {\Large\bfseries MA1001B - Statistical Modeling for Decision Making\par}
  \vspace{0.7cm}
  {\large Lesson 24: Central Limit Theorem\par}
  \vspace{0.7cm}
  {\color{TecRojo}\rule{\textwidth}{1.2pt}}
  \vspace{0.8cm}

  MA1001B Faculty\\[0.3cm]
  Term\\[0.3cm]
  Tecnol\'ogico de Monterrey
\end{frame}

\begin{frame}{The Skewed-Population Question}
  \begin{alertblock}{Guiding question}
    If delay times are exponential rather than bell-shaped, when is the
    sampling distribution of $\bar x$ still approximately normal?
  \end{alertblock}

  \vspace{0.6em}
  By the end of this lesson, you can:
  \begin{itemize}
    \item recognize a right-skewed population that is not normal;
    \item compare $\bar x$ histograms for $n=4$ and $n=40$;
    \item state the CLT for the sample mean;
    \item explain why a two-SE tail fails at small $n$.
  \end{itemize}

  \vfill
  \scriptsize All delay times used today are synthetic. No real company data are used.
\end{frame}

\begin{frame}{An Exponential Delay Clock}
  \small
  $X\sim\mathrm{Exponential}(\text{mean}=10)$:
  \[
    E[X]=\sigma=10,\qquad \text{skewness}=2,\qquad \text{median}=6.931472.
  \]
  The right tail is long: $P(X>20)=0.135335$.
  For iid samples,
  \[
    E[\bar x]=10,\qquad \mathrm{SE}(\bar x)=10/\sqrt{n}.
  \]

  \vspace{0.3em}
  \begin{exampleblock}{CLT target}
    The CLT is a statement about $\bar x$, not about the original delays
    becoming normal.
  \end{exampleblock}
\end{frame}

\begin{frame}[fragile]{Step 1: Skewed Population}
  \begin{columns}[T]
    \begin{column}{0.42\textwidth}
      \small
      \begin{lstlisting}
model = expon(scale=10)
mean = model.mean()
skew = model.stats(
  moments="s"
)
      \end{lstlisting}

      \begin{block}{Verified output}
        Mean $=$ sd $=10$\\
        Median $=6.931472$\\
        Skewness $=2.000000$
      \end{block}
    \end{column}
    \begin{column}{0.58\textwidth}
      \centering
      \includegraphics[width=\textwidth]{../figures/clt_01_skewed_population.png}
      \vspace{0.2em}
      \scriptsize Exponential delay density from Step 1
    \end{column}
  \end{columns}
\end{frame}

\begin{frame}{Larger $n$ Pulls $\bar x$ Toward a Bell}
  \small
  5{,}000 seed-42 samples:
  \[
    n=4:\ \widehat{\mathrm{SE}}=5.076401,\ \widehat{\mathrm{skew}}=1.049897.
  \]
  \[
    n=40:\ \widehat{\mathrm{SE}}=1.574371,\ \widehat{\mathrm{skew}}=0.332572.
  \]
  Theoretical skew of an exponential mean is $2/\sqrt{n}$. That value is
  $1$ at $n=4$ and $0.316228$ at $n=40$.
\end{frame}

\begin{frame}[fragile]{Step 2: $n=4$ versus $n=40$}
  \begin{columns}[T]
    \begin{column}{0.40\textwidth}
      \small
      \begin{lstlisting}
draws = model.rvs(
  size=(5000, n)
)
xbars = draws.mean(axis=1)
      \end{lstlisting}

      \begin{block}{Verified output}
        $n=4$ skew $=1.049897$\\
        $n=40$ skew $=0.332572$
      \end{block}
    \end{column}
    \begin{column}{0.60\textwidth}
      \centering
      \includegraphics[width=\textwidth]{../figures/clt_02_means_n4_n40.png}
      \vspace{0.2em}
      \scriptsize Sampling histograms from Step 2
    \end{column}
  \end{columns}
\end{frame}

\begin{frame}{Step 3: Small $n$ Makes the Tail Wrong}
  \begin{columns}[T]
    \begin{column}{0.42\textwidth}
      \small
      Exact law: $\bar x\sim\mathrm{Gamma}$.
      Event: $\bar x>\mu+2\,\mathrm{SE}$.

      \vspace{0.4em}
      $n=4$: exact $0.042380$ versus normal $0.022750$.\\
      $n=40$: exact $0.030560$ versus normal $0.022750$.

      \vspace{0.4em}
      \textbf{Limit:} small $n$ plus skew makes the CLT a poor
      approximation. Lesson 25 studies $\hat p$.
    \end{column}
    \begin{column}{0.58\textwidth}
      \centering
      \includegraphics[width=\textwidth]{../figures/clt_03_small_n_skew_limit.png}
      \vspace{0.2em}
      \scriptsize Exact versus normal two-SE tails from Step 3
    \end{column}
  \end{columns}
\end{frame}

\begin{frame}{Concept Check}
  \begin{enumerate}
    \item Why is the delay population not a normal model?
    \item What happens to the skewness of $\bar x$ as $n$ grows?
    \item At $n=4$, compare the exact two-SE tail with $0.022750$.
    \item Does $n=40$ make the normal tail exact?
  \end{enumerate}

  \vspace{0.6em}
  \begin{block}{Connection to Lesson 25}
    The session can continue directly into the sampling distribution of a
    sample proportion $\hat p$.
  \end{block}

  \vfill
  \scriptsize Source: OpenStax, \textit{Introductory Business Statistics 2e},
  Chapter 7, Sections 7.1--7.2. CC BY-NC-SA.
\end{frame}

\appendix



\end{document}
